\begin{frame}{Medidas de gasto}
    \begin{itemize}
        \item \textbf{FEPC: } Eliminar subsidio al diésel.
        \item \textbf{Inversión:} 
        \begin{itemize}
            \item Mantener la suspensión de vigencias futuras.
            \item Llevar al mínimo el gasto de inversión discrecional.
            \item Flexibilizar APP.
        \end{itemize} 
        \item \textbf{Personal:} Revertir el aumento presupuestado del gasto de personal.
        \item \textbf{Salud:} 
        \begin{itemize}
            \item Acotar y condicionar los beneficios del PBS. 
            %\item El saneamiento de las finanzas del sector salud no puede darse a costa de las finanzas del GNC.
        \end{itemize} 
    \end{itemize}
\end{frame}

\begin{frame}{Medidas de gasto}
    \begin{itemize}
            \item \textbf{Pensiones:} 
            \begin{itemize}
                \item Reforma Paramétrica. 
                \item Impuesto a Pensiones, \item Contribuciones al Fondo de Solidaridad Pensional.
            \end{itemize}  
        \item \textbf{SGP:} 
        \begin{itemize}
            \item Aplazar la Ley de Competencias hasta que no culmine la consolidación fiscal.
            \item Preferiblemente introducir un nuevo Acto Legislativo para definir SGP como \% del PIB.
        \end{itemize}    
    \end{itemize}
\end{frame}
